\documentclass[11pt,reqno]{amsart}
\usepackage[margin=1in]{geometry}
\usepackage[expansion=false]{microtype}
\usepackage{booktabs}
\usepackage{xcolor}
\usepackage[colorlinks=true,linkcolor=blue!55!black,citecolor=blue!55!black,urlcolor=blue!55!black]{hyperref}
\newtheorem{theorem}{Theorem}
\newtheorem{lemma}[theorem]{Lemma}
\newtheorem{proposition}[theorem]{Proposition}
\theoremstyle{remark}
\newtheorem{remark}[theorem]{Remark}
\newcommand{\Q}{\mathbb Q}
\newcommand{\Z}{\mathbb Z}
\newcommand{\F}{\mathbb F}
\newcommand{\Jac}{\operatorname{Jac}}
\newcommand{\Sel}{\operatorname{Sel}}
\DeclareMathOperator{\Sha}{Sha}
\DeclareMathOperator{\divisor}{div}
\hypersetup{pdftitle={A genus-three hyperelliptic curve with rational Jacobian three-torsion of rank three},pdfauthor={Roy van Rijn}}
\title[A genus-three curve with rational three-torsion]{A genus-three hyperelliptic curve with rational Jacobian three-torsion of rank three}
\author{Roy van Rijn}
\subjclass[2020]{11G30, 11G10, 14H40}
\keywords{hyperelliptic curve, rational torsion, split Jacobian, Selmer group}
\date{}
\begin{document}
\begin{abstract}
We exhibit a hyperelliptic curve of genus three over the rationals whose
Jacobian has rational three-torsion isomorphic to \((\Z/3\Z)^3\).
This supplies the genus-three example asked for by Levin, Yan and Wiljanen.
The curve is a double cover of a rational conic with two commuting sign
involutions. Three explicit maps of degree eight to elliptic curves with
rational points of order three produce independent torsion classes; the
Weil pairing gives the matching upper bound. Pullback and pushforward
give dual isogenies of degree \(2^9\). Consequently, for every quadratic
twist and every \(n\geq1\), its \(3^n\)-Selmer group and
\(3^n\)-torsion in the Tate--Shafarevich group identify with those of
the product of the three elliptic twists, without a finiteness assumption.
We also determine all eight rational points and the full Mordell--Weil
group, which is finite with invariant factors \([3,6,12]\).
\end{abstract}
\maketitle

\section{The existence question and the example}

For a smooth projective curve \(C/\Q\) of genus \(g\), the Weil pairing
bounds the dimension of \(\Jac(C)(\Q)[3]\) over \(\F_3\) by \(g\).
Levin, Yan and Wiljanen \cite[Theorem~9 and the following discussion]{LYW}
give hyperelliptic examples attaining this bound in genera one, two and
four. Their rank-three example has genus four, and they ask whether
a hyperelliptic curve with genus three and rational Jacobian
three-torsion of rank three exists. We give such a curve.

\begin{theorem}\label{thm:main}
Let \(C/\Q\) be the smooth projective model of \(R^2=F(t)\), where
\begin{equation}\label{eq:octic}
\begin{split}
 F(t)={}&t^8-6t^7-t^6+18t^5+93t^4\\
        &{}-36t^3-4t^2+48t+16.
\end{split}
\end{equation}
Then \(C\) is hyperelliptic of genus three and
\[
 \Jac(C)(\Q)[3]\simeq(\Z/3\Z)^3.
\]
For \((h_1,h_2,h_3)=(3/2,-3/4,-7/6)\), set
\begin{equation}\label{eq:elliptic}
 E_i:\quad Y^2=X^3+(h_iX+h_i+1)^2.
\end{equation}
There are explicit morphisms \(\pi_i:C\to E_i\) of degree eight.
Writing \(A=E_1\times E_2\times E_3\) and \(J=\Jac(C)\), the maps
\[
 \Phi=\sum_i\pi_i^*:A\longrightarrow J,
 \qquad
 \Psi=(\pi_{1*},\pi_{2*},\pi_{3*}):J\longrightarrow A
\]
are dual isogenies satisfying
\begin{equation}\label{eq:composites}
 \Psi\Phi=[8]_A,\qquad \Phi\Psi=[8]_J,
 \qquad \deg\Phi=\deg\Psi=2^9.
\end{equation}
\end{theorem}

The construction belongs to the classical method of splitting a
Jacobian into elliptic factors. In particular,
Howe, Lepr\'evost and Poonen \cite[Proposition~14]{HLP} construct
hyperelliptic genus-three Jacobians by gluing three elliptic curves
when their twisting factor vanishes. Our proof gives the conic,
the three quotients and all maps directly; it does not require
an application of that proposition with implicit gluing data.

The reason three independent classes survive is simple. Each elliptic
curve in \eqref{eq:elliptic} has a rational point of order three.
Pulling it back along a degree-eight map cannot kill it, because
pushforward followed by pullback is multiplication by eight on the
elliptic curve. Independence is a different issue: three maps alone
would not suffice. Here the two sign involutions place the three
pullbacks in different characters of a group of order four. Since four
is invertible modulo three, these characters separate the classes.
The same separation makes the off-diagonal push--pull maps vanish,
and gives \eqref{eq:composites}.

The assertion concerns \(J(\Q)[3]\), the subgroup killed by three.
Theorem~\ref{thm:arithmetic} below determines the full rational torsion.
The Jacobian is geometrically nonsimple because of its elliptic factors.
M\"uller and Reitsma \cite[Theorem~1.1 and Remark~1.2]{MR} distinguish
the occurrence of an exact full torsion group from that of a subgroup,
and study geometrically simple examples. Theorem~\ref{thm:main}
addresses the existence question above. The full torsion in our example
has invariant factors \([3,6,12]\), so it does not realize their possible
exception \([3,3,3]\), even without the geometric simplicity condition.

The isogeny has a further consequence. For an abelian variety \(B/\Q\)
and a squareclass \(d\in\Q^\times/\Q^{\times2}\), let \(B^d\) denote
the quadratic twist by the automorphism \([-1]\). We use
\(\Sel_m(B/\Q)\) for the Selmer group for multiplication by \(m\),
and \(\Sha(B/\Q)\) for the Tate--Shafarevich group.

\begin{theorem}\label{thm:descent}
For the varieties in Theorem~\ref{thm:main}, every rational
squareclass \(d\) and every integer \(n\geq1\), the twisted maps induce
isomorphisms
\begin{align*}
 \Sel_{3^n}(J^d/\Q)&\simeq
       \bigoplus_{i=1}^3\Sel_{3^n}(E_i^d/\Q),\\
 \Sha(J^d/\Q)[3^n]&\simeq
       \bigoplus_{i=1}^3\Sha(E_i^d/\Q)[3^n].
\end{align*}
They also identify the full three-primary Tate--Shafarevich groups.
No finiteness assumption on these groups is required.
\end{theorem}

\begin{theorem}\label{thm:arithmetic}
For the curve of Theorem~\ref{thm:main},
\[
 C(\Q)=\{(-1,\pm9),(0,\pm4),(2,\pm36),\infty_+,\infty_-\},
 \qquad J(\Q)\simeq\Z/3\Z\oplus\Z/6\Z\oplus\Z/12\Z.
\]
Here \(R/t^4\) has value \(\pm1\) at \(\infty_\pm\).
In particular \(J\) has rank zero and 216 rational points.
Its conductor is \(243000\).
\end{theorem}

\section{The conic and its three elliptic quotients}

Put \(d_0=3(t^2+2)\) and define functions on \(C\) by
\begin{equation}\label{eq:parameter}
 x=\frac{t^2-4t-2}{d_0},\qquad
 y=\frac{-2t^2-4t+4}{d_0},\qquad
 w=\frac{12R}{d_0^2}.
\end{equation}
They satisfy
\begin{equation}\label{eq:conic}
 6x^2+3y^2=2,\qquad
 w^2=-\frac92x^4-\frac{117}{16}y^4+\frac{59}{18}.
\end{equation}
Conversely, generically \(t=(y-2/3)/(x-1/3)\). Thus
\eqref{eq:conic} is a birational model of \(C\); it comes from
parametrizing the conic by lines through \((1/3,2/3)\).

Modulo seven the polynomial \(F\) and its derivative are coprime.
As \(F\) is monic of degree eight, it is squarefree over \(\Q\).
The double cover of \(\mathbf P^1\) therefore has eight simple branch
points, and Riemann--Hurwitz gives genus three. In particular all
nonconstant rational maps below extend uniquely to morphisms of
smooth projective curves. Apparent poles in their formulas create no
additional restrictions on the theorem.

The conic model has commuting involutions
\[
 \sigma_x(x,y,w)=(-x,y,w),\qquad
 \sigma_y(x,y,w)=(x,-y,w).
\]
Together they generate a group \(G\simeq(\Z/2\Z)^2\).
The quotients by \(\sigma_x,\sigma_y,\sigma_x\sigma_y\), respectively,
have coordinates and equations
\begin{equation}\label{eq:rows}
\begin{array}{c|cc|rrr}
 i&z_i&v_i&a_i&b_i&c_i\\ \hline
 1&y&w&-135/16&3/2&25/9\\
 2&x&w/2&-135/16&39/8&1/144\\
 3&x/y&-w/(3y^2)&25/9&59/18&1/144
\end{array}
\qquad
 v_i^2=a_i z_i^4+b_i z_i^2+c_i.
\end{equation}
For the first two rows, eliminate \(x^2\) or \(y^2\) in
\eqref{eq:conic}. For the third, use
\(y^2=2/(3(2z_3^2+1))\). These substitutions give the displayed
equations. They also show that each fixed field is generated by
the stated \(z_i,v_i\): adjoining \(x\), \(y\), or \(y\), respectively,
recovers \(\Q(C)\) in degree two.

Each quartic on the right in \eqref{eq:rows} is squarefree:
\(a_ic_i\ne0\) and \(b_i^2-4a_ic_i\ne0\).
Its smooth projective double cover has genus one, with a rational
point \((0,\sqrt{c_i})\). We next give the degree-four map from
each such curve to the desired elliptic target.

\begin{lemma}\label{lem:quartic}
Let \(a,b,c\in\Q\) with \(ac(b^2-4ac)\ne0\).
On the genus-one curve \(v^2=az^4+bz^2+c\), set
\begin{equation}\label{eq:quartic-map}
 U=\frac{v^2}{4z^2},\qquad
 V=\frac{v}{8z}\left(az^2-\frac{c}{z^2}\right).
\end{equation}
These functions define a morphism of degree four to
\begin{equation}\label{eq:target}
 V^2=U\left(U^2-\frac b2U+\frac{b^2-4ac}{16}\right).
\end{equation}
The involution \((z,v)\mapsto(-z,v)\) sends \((U,V)\) to \((U,-V)\).
\end{lemma}
\begin{proof}
Directly,
\[
 \frac{V^2}{U}=\frac1{16}(az^2-c/z^2)^2
             =U^2-\frac b2U+\frac{b^2-4ac}{16}.
\]
The cubic in \eqref{eq:target} is squarefree by the hypotheses:
its quadratic factor has nonzero constant term and discriminant
\(ac\). Hence the target is elliptic. As a function of \(z\),
\[
 U=\frac{az^4+bz^2+c}{4z^2}
\]
has degree four. The map \(z\) from the quartic curve to
\(\mathbf P^1\) has degree two, while \(U\) on the target elliptic
curve has degree two. The tower of function fields therefore gives
degree \((2\cdot4)/2=4\) for the displayed map.
The final assertion follows directly from the formulas.
\end{proof}

Apply Lemma~\ref{lem:quartic} to each row, and put
\(X_i=U_i-1\), \(Y_i=V_i\).
In all three cases substitution into \eqref{eq:target} gives
\eqref{eq:elliptic}. More explicitly the coefficient identities are
\[
 h_i^2=3-b_i/2,\quad
 2h_i(h_i+1)=3-b_i+(b_i^2-4a_ic_i)/16,
 \quad
 (h_i+1)^2=1-b_i/2+(b_i^2-4a_ic_i)/16.
\]
Composing the quotient map of degree two with this map of degree
four gives \(\pi_i\) of degree eight. The equations of the targets
are nonsingular, since their discriminants are
\(16(h_i+1)^3(4h_i^3-27(h_i+1))\ne0\).

The action on the maps is
\begin{equation}\label{eq:characters}
\begin{array}{c|rr}
 &\sigma_x&\sigma_y\\ \hline
 \pi_1&+&-\\
 \pi_2&-&+\\
 \pi_3&-&-
\end{array}
\end{equation}
where a minus sign means composition with \([-1]\) on the elliptic
curve. For \(\pi_3\), either sign change sends \(z_3\) to
\(-z_3\) and fixes \(v_3\); for the other two maps this is immediate
from \eqref{eq:rows}. The lemma proves all the signs.

\section{The three rational torsion classes}

On \(E_i\), let \(O_i\) be the point at infinity and put
\[
 P_i=(0,-h_i-1),\qquad f_i=Y+h_iX+h_i+1.
\]
The line \(f_i=0\) meets the cubic only at \(P_i\), with
intersection multiplicity three, because substitution leaves
\(X^3=0\). Also \(f_i\) has a pole of order three at \(O_i\).
Consequently
\begin{equation}\label{eq:divisor}
 \divisor(f_i)=3(P_i)-3(O_i).
\end{equation}
Since \(h_i+1\ne0\), the point \(P_i\ne O_i\) has order exactly three.

Define \(D_i=\pi_i^*([P_i-O_i])\in J(\Q)[3]\).
Push--pull for finite morphisms of curves gives
\[
 \pi_{i*}\pi_i^*=[\deg\pi_i]=[8]
 \quad\hbox{on }E_i=\Jac(E_i).
\]
It follows that \(D_i\ne0\). By \eqref{eq:characters}, the action
of \(G\) on \(D_i\) is the corresponding nontrivial character
\(\chi_i:G\to\{\pm1\}\).
For clarity, if \(\sum_i e_iD_i=0\) with \(e_i\in\F_3\), apply
\[
 \frac14\sum_{g\in G}\chi_j(g)g^*
\]
on \(J[3]\). Character orthogonality yields \(e_jD_j=0\).
Thus each \(e_j=0\), and the three classes are independent.

Finally the canonical principal polarization gives a nondegenerate
Galois-equivariant alternating pairing
\[
 e_3:J[3]\times J[3]\longrightarrow\mu_3.
\]
On rational points its value lies in \(\mu_3(\Q)=\{1\}\).
Hence \(J(\Q)[3]\) is an isotropic subspace of the six-dimensional
symplectic \(\F_3\)-space \(J[3](\overline\Q)\), so its dimension
is at most three. This proves the torsion assertion of
Theorem~\ref{thm:main}.

The pulled-back functions \(f_i\circ\pi_i\) make the classes explicit
as cube-divisor functions by \eqref{eq:divisor}.
The proof uses their divisors on the curve and requires no
specialization to quadratic number fields.

\section{Dual isogenies and descent on every quadratic twist}

We first finish the isogeny assertion in Theorem~\ref{thm:main}.
If \(i\ne j\), choose \(g\in G\) with
\(\chi_i(g)=-\chi_j(g)\). Its action \(g^*\) on \(J\) is an
involution. Equations \eqref{eq:characters} give
\[
 g^*\pi_j^*=\chi_j(g)\pi_j^*,\qquad
 \pi_{i*}g^*=\chi_i(g)\pi_{i*}.
\]
Therefore \(\pi_{i*}\pi_j^*=-\pi_{i*}\pi_j^*\).
The group of homomorphisms of abelian varieties has no integer
torsion: the image of a homomorphism killed by a nonzero integer
is connected and contained in a finite group scheme, and in
characteristic zero must be zero. Thus all cross terms vanish.
The diagonal terms are \([8]\), proving \(\Psi\Phi=[8]_A\).

It follows that \(\ker\Phi\subset A[8]\) is finite. Since \(A\)
and \(J\) both have dimension three, \(\Phi\) is an isogeny.
Now \((\Phi\Psi-[8]_J)\Phi=0\), and surjectivity of \(\Phi\)
gives \(\Phi\Psi=[8]_J\). Pullback and pushforward are dual
under the canonical principal polarizations of curve Jacobians.
Hence \(\deg\Phi=\deg\Psi\), and
\[
 (\deg\Phi)^2=\deg[8]_A=8^6
\]
gives \(\deg\Phi=8^3=2^9\), as claimed.

\begin{proof}[Proof of Theorem~\ref{thm:descent}]
The homomorphisms \(\Phi\) and \(\Psi\) are defined over \(\Q\)
and commute with \([-1]\). They therefore descend on every
quadratic twist to maps \(\Phi^d:A^d\to J^d\) and
\(\Psi^d:J^d\to A^d\) with the same composites \([8]\).
Here \(A^d=E_1^d\times E_2^d\times E_3^d\).

Fix \(m=3^n\), and choose \(a\in\Z\) with \(8a\equiv1\pmod m\).
On \(m\)-torsion, \(\Phi^d\) and \(a\Psi^d\) are inverse
Galois-module maps. They induce inverse maps on global and
local first cohomology of these torsion modules.
For each completion \(k=\Q_v\), they also give inverse maps
\[
 A^d(k)/mA^d(k)\ \longleftrightarrow\ J^d(k)/mJ^d(k).
\]
Indeed the composites are multiplication by \(8a\), which is the
identity modulo \(m\). Naturality in the Kummer exact sequence
\[
 0\longrightarrow B[m]\longrightarrow B
 \xrightarrow{[m]}B\longrightarrow0
\]
identifies the images of these quotients in \(H^1(k,B[m])\).
This argument applies at every place, including three and the real
place. Defining the Selmer group as the global cohomology classes
lying in all these local Kummer images proves the first isomorphism.
Cohomology and the Kummer conditions for a finite product are the
direct sums of those for its factors.

For the second assertion use
\[
 \Sha(B/\Q)=\ker\left(H^1(\Q,B)\longrightarrow
                              \prod_v H^1(\Q_v,B)\right).
\]
Functoriality gives maps on these kernels, with composites \([8]\).
On the subgroups killed by \(3^n\), multiplication by \(a\)
again gives inverse maps. On the union of these subgroups,
multiplication by eight is an automorphism: on every element
of finite three-power order its inverse is the unique inverse
modulo that order. Thus the full three-primary groups are also
isomorphic. The argument uses no finiteness assumption.
\end{proof}

\section{Rational points and full rational torsion}

\begin{proof}[Proof of Theorem~\ref{thm:arithmetic}]
Exact elliptic two-descent gives proved rank bounds \((0,0)\) for each
of the three curves \eqref{eq:elliptic}; these computations, with the
rank database disabled, are included in the accompanying SageMath code.
The respective Cremona labels are \(90a1,90b1,30a2\), with conductors
\(90,90,30\). Thus the isogeny gives rank zero and conductor
\(90\cdot90\cdot30=243000\) for \(J\).

First determine the rational points. In the coordinates of
\eqref{eq:elliptic}, the six points of \(E_1(\Q)\) are
\[
 O,\quad(-1,0),\quad(0,\pm5/2),\quad(5,\pm15).
\]
There is no rational point at infinity on the conic
\(6x^2+3y^2=2\), and \(y=0\) would require \(x^2=1/3\).
Consequently \(U=w^2/(4y^2)=X+1\) is finite and belongs to
\(\{0,1,6\}\). Its fibre satisfies
\[
 -\frac{135}{16}y^4+\left(\frac32-4U\right)y^2+\frac{25}{9}=0.
\]
For \(U=0,6\), the discriminants of these quadratics in \(y^2\)
are 96 and 600 respectively, both nonsquares in \(\Q\).
For \(U=1\) the polynomial factors as
\[
 -\frac{135}{16}(y-2/3)(y+2/3)(y^2+20/27).
\]
It follows that \(y=\pm2/3\), then \(x=\pm1/3\) and
\(w=\pm4/3\). The parametrization \eqref{eq:parameter} gives
\(t=-1,0,2,\infty\). Substitution into \eqref{eq:octic} gives
exactly the eight stated points, including both branches at infinity.

To determine torsion, good reduction and the elliptic decomposition give
\[
\begin{array}{c|rrrrr}
 p&7&11&13&17&19\\ \hline
 \#J(\F_p)&432&1296&3888&3456&13824.
\end{array}
\]
The gcd is \(432=16\cdot27\). Thus Theorem~\ref{thm:main}
already determines the full three-primary subgroup, and it remains
to determine the two-primary subgroup \(H\).

We use the following small finite divisor-class computation at seven.
Over \(\F_7\), put
\[
 I=[\infty_- -\infty_+],\quad P=[(0,4)-\infty_+],\quad
 Q=[D-2\infty_+],\quad D:\ (t^2+1,R-t+2)=0.
\]
Here \(D\) is a closed point of degree two. The three classes
\(27I,27P,27Q\) are killed by four and generate groups of orders
4, 8, 16 successively; precisely four of the resulting sixteen classes
are killed by two. This is an exact finite certificate, verified by
divisor arithmetic in the supplied code. Since the two-part of
\(\#J(\F_7)\) is 16, it proves
\[
 J(\F_7)[2^\infty]\simeq(\Z/4\Z)^2,
 \qquad H\ \lhook\joinrel\longrightarrow\ (\Z/4\Z)^2.
\]
The same certificate shows that the reduction of \(3I\) has order four.

We next exhibit a rational point of order two that is not a double.
Let \(T=3[(0,4)-\infty_+]\), and set
\[
 g=t^6-3t^5-3t^4-12t^3+6t^2-12t-8.
\]
The identity \(g^2-F(t)(t^2+2)^2=-200t^6\) gives
\[
 \divisor\bigl(g+(t^2+2)R\bigr)=6(0,4)-6\infty_+,
\]
so \(2T=0\). To show that \(T\) is not a double, work on the good
reduction at 17. For \(L=\F_{17}[\theta]/(F(\theta))\), the
even-degree descent homomorphism
\[
 J(\F_{17})/2J(\F_{17})\longrightarrow
                  L^\times/(L^{\times2}\F_{17}^\times)
\]
sends \([(a,b)-\infty_+]\) to \(a-\theta\).
The rational infinity has trivial image; see
Poonen--Schaefer \cite[Section~5, Proposition~5.1 and (8)]{PS}.
The rational point identifies degree-zero rational divisor classes with
the Jacobian points, as required here. We use only that doubles have
trivial image, not injectivity of the fake descent map.
The image of \(T\) is \(-\theta\). The polynomial \(t^2+5\) is an
irreducible factor of \(F\bmod17\); in the corresponding factor
\(\F_{17^2}\), the norm of \(-\theta\) is 5. This is a nonsquare in
\(\F_{17}\), and multiplying by a scalar \(c\ne0\) changes the norm
to \(5c^2\). Hence the image is nontrivial, even modulo scalars.
Thus \(T\ne0\) and \(T\notin2J(\Q)\).

Finally an explicit Pell identity supplies rational four-torsion.
Put
\[
\begin{split}
 A={}&-t^8+8t^7-4t^6-52t^5-115t^4\\
     &{}+296t^3+746t^2+764t-139.
\end{split}
\]
Let \(B\) be the polynomial part of \(A\sqrt F\) at infinity, with
\(\sqrt F=t^4-3t^3-5t^2+\cdots\). Exact multiplication gives
\(B^2-FA^2=19440000\); the code includes both expanded polynomials.
The leading terms give
\[
 \divisor(B+AR)=12\infty_- -12\infty_+.
\]
Thus \(3I\) is killed by four, and its reduction at seven shows its
order is exactly four. Its nonzero double cannot equal \(T\), since
\(T\) is not divisible by two. Consequently
\(\Z/4\Z\oplus\Z/2\Z\subseteq H\).
A larger subgroup of \((\Z/4\Z)^2\) would be that whole group,
in which every element of order two is a double, contradicting the
property of \(T\). Therefore \(H\simeq\Z/4\Z\oplus\Z/2\Z\).
Combining the two primary parts proves the asserted invariant factors.
\end{proof}

\paragraph{Verification.}
All defining identities and maps are given above over \(\Q\).
An accompanying SageMath checker verifies squarefreeness modulo seven,
the conic and quotient identities, the degrees of the three maps,
the orders of the marked elliptic points, and the sign characters.
These exact algebraic checks support the displayed construction.
The independence, Weil-pairing bound, duality and compatibility with
all local Kummer conditions are the written arguments above;
the checker does not compute Selmer or Tate--Shafarevich groups.
The arithmetic checker also computes proved elliptic rank bounds,
the target torsion points, all rational fibres, the two principal-divisor
certificates, the complete sixteen-element Sylow subgroup at seven and
the square-class obstruction at 17. These calculations use exact
rational and finite-field arithmetic in SageMath 10.9, with no height
search or random input.

\paragraph{AI assistance and review.}
ChatGPT (OpenAI) assisted with mathematical arguments, literature checks,
exposition, proof review and verification code.

\begin{thebibliography}{9}
\bibitem{HLP}
E.~W.~Howe, F.~Lepr\'evost and B.~Poonen,
\emph{Large torsion subgroups of split Jacobians of curves of genus two or three},
Forum Math. \textbf{12} (2000), 315--364.
\href{https://arxiv.org/abs/math/9809210}{arXiv:math/9809210}.

\bibitem{LYW}
A.~Levin, S.~Yan and L.~Wiljanen,
\emph{Quadratic fields with a class group of large \(3\)-rank},
\href{https://arxiv.org/abs/1910.12276v1}{arXiv:1910.12276v1} (2019).

\bibitem{MR}
J.~S.~M\"uller and B.~Reitsma,
\emph{Computing torsion subgroups of Jacobians of hyperelliptic curves of genus \(3\)},
\href{https://arxiv.org/abs/2211.03372v4}{arXiv:2211.03372v4} (2023).

\bibitem{PS}
B.~Poonen and E.~F.~Schaefer,
\emph{Explicit descent for Jacobians of cyclic covers of the projective line},
J. reine angew. Math. \textbf{488} (1997), 141--188.
\href{https://math.mit.edu/~poonen/papers/descent.pdf}{Author's version}.
\end{thebibliography}
\end{document}
